\documentclass[14pt,a4paper,oneside]{extbook}
\usepackage[utf8]{vietnam}
\usepackage{amsmath,amssymb,mathrsfs,fancyhdr,enumerate,multirow,makecell}
\usepackage{tkz-euclide,tikz-3dplot,tikz,tkz-tab}
\usetikzlibrary{arrows.meta,calc,intersections,angles,quotes,backgrounds,shapes.geometric,patterns,through,shadows,positioning}
\usepackage{yhmath}
\usepackage[top=0.7cm, bottom=0.7cm, left=1.0cm, right=1.0cm, headheight=16pt, headsep=3mm, footskip=5mm]{geometry}
\usepackage[hidelinks,unicode]{hyperref}
\usepackage{qrcode}
\usepackage[many]{tcolorbox}
\usepackage{fontawesome}

\pagestyle{fancy}
\fancyhf{}
\fancyhead[L]{\footnotesize\bfseries\color{teal!85!black}\faBook~SỔ TAY KIẾN THỨC TOÁN 5}
\fancyhead[R]{\footnotesize\bfseries\color{blue!85!black}Lớp Toán Thầy Lĩnh -- loptoanthaylinh.vn}
\fancyfoot[C]{\footnotesize\bfseries\color{gray!80!black}-- Sổ tay kiến thức: Bài 37. Tỉ lệ bản đồ và ứng dụng --}
\renewcommand{\headrulewidth}{0.4pt}

\begin{document}

% --- BANNER TIÊU ĐỀ BÀI HỌC ---
\noindent
\begin{tcolorbox}[
  enhanced,
  colback=teal!8!white,
  colframe=teal!80!black,
  arc=3mm,
  boxrule=1.1pt,
  width=\linewidth,
  halign=center,
  top=2.2mm, bottom=2.2mm,
  left=3mm, right=3mm,
  drop fuzzy shadow
]
  {\large\bfseries\color{teal!95!black}\faBookmark\ BÀI 37. TỈ LỆ BẢN ĐỒ VÀ ỨNG DỤNG}\\[2pt]
  {\scriptsize\color{black!75}\itshape Ý nghĩa tỉ lệ 1:n, công thức tính khoảng cách thực tế và khoảng cách thu nhỏ trên 1 trang A4}
\end{tcolorbox}

\vspace{3.2mm}
% --- KHỐI 1: KHÁI NIỆM TỈ LỆ BẢN ĐỒ ---
\noindent
\begin{minipage}[t]{0.56\textwidth}
  \begin{tcolorbox}[
    enhanced,
    title={\footnotesize\bfseries\color{white}\faMapO\ 1. Khái niệm Tỉ lệ bản đồ},
    colback=blue!2!white,
    colframe=blue!75!black,
    arc=2.5mm, boxrule=0.8pt, top=1.8mm, bottom=1.8mm, left=2.2mm, right=2.2mm
  ]
    \footnotesize
    $\bullet$ Tỉ lệ bản đồ được ghi dưới dạng: \textbf{$1 : n$} hoặc $\dfrac{1}{n}$.\\
    $\bullet$ \textbf{Ý nghĩa}: Khoảng cách $1$ đơn vị trên bản đồ tương ứng với \textbf{$n$ đơn vị đó ngoài thực tế}.\\
    $\triangleright$ \textit{VD: Bản đồ tỉ lệ $1 : 1\,000$ nghĩa là:}\\
    \quad $1\text{ cm}$ trên bản đồ $= 1\,000\text{ cm} = 10\text{ m}$ ngoài thực tế.\\
    $\triangleright$ \textit{VD: Bản đồ tỉ lệ $1 : 1\,000\,000$ nghĩa là:}\\
    \quad $1\text{ cm}$ trên bản đồ $= 1\,000\,000\text{ cm} = 10\text{ km}$ ngoài thực tế.
  \end{tcolorbox}
\end{minipage}%
\hfill
\begin{minipage}[t]{0.42\textwidth}
  \begin{tcolorbox}[
    enhanced,
    title={\footnotesize\bfseries\color{white}\faSearch\ 2. Nguyên tắc đồng vị đo},
    colback=teal!2!white,
    colframe=teal!75!black,
    arc=2.5mm, boxrule=0.8pt, top=1.8mm, bottom=1.8mm, left=2mm, right=2mm
  ]
    \footnotesize
    $\bullet$ Trên bản đồ và ngoài thực tế khi lập tỉ số \textbf{phải cùng đơn vị đo}.\\
    $\bullet$ Đổi nhanh:\\
    $1\text{ km} = 1\,000\text{ m} = 100\,000\text{ cm}$.\\
    $1\text{ m} = 100\text{ cm} = 1\,000\text{ mm}$.
  \end{tcolorbox}
\end{minipage}

\vspace{3mm}

% --- KHỐI 2: HAI BÀI TOÁN CƠ BẢN VỀ TỈ LỆ BẢN ĐỒ ---
\noindent
\begin{tcolorbox}[
  enhanced,
  title={\footnotesize\bfseries\color{white}\faCalculator\ 3. Hai bài toán ứng dụng cơ bản},
  colback=purple!2!white,
  colframe=purple!70!black,
  arc=2.5mm, boxrule=0.8pt, top=1.8mm, bottom=1.8mm, left=2.5mm, right=2.5mm
]
  \footnotesize
  \begin{tabular}{p{0.48\textwidth}|p{0.48\textwidth}}
    \textbf{Dạng 1: TÌM ĐỘ DÀI THỰC TẾ} & \textbf{Dạng 2: TÌM ĐỘ DÀI THU NHỎ (TRÊN BẢN ĐỒ)} \\ \hline
    Biết độ dài trên bản đồ $a$, tỉ lệ $1 : n$: & Biết độ dài thực tế $A$, tỉ lệ $1 : n$: \\
    \centering\textbf{\color{purple!90!black}$\text{Độ dài thực tế} = a \times n$} & \centering\textbf{\color{blue!90!black}$\text{Độ dài trên bản đồ} = A : n$} \\[2pt]
    \textit{(Sau đó đổi ra m hoặc km theo yêu cầu)} & \textit{(Trước tiên đổi A về cùng đơn vị với bản đồ: cm)}
  \end{tabular}
\end{tcolorbox}

\vspace{3mm}

% --- KHỐI 3: SƠ ĐỒ TƯ DUY ---
\noindent
\begin{tcolorbox}[
  enhanced,
  title={\footnotesize\bfseries\color{white}\faSitemap\ Sơ đồ tư duy: Xử lý bài toán tỉ lệ bản đồ},
  colback=yellow!2!white,
  colframe=teal!85!black,
  arc=3mm, boxrule=0.9pt, top=1.8mm, bottom=1.8mm, left=1.5mm, right=1.5mm,
  halign=center
]
  \begin{tikzpicture}[
    scale=0.92,
    every node/.style={transform shape},
    b1/.style={rectangle, rounded corners=4pt, fill=blue!15, draw=blue!70!black, text=blue!90!black, font=\bfseries\scriptsize, inner sep=4pt},
    b2/.style={rectangle, rounded corners=4pt, fill=orange!15, draw=orange!70!black, text=orange!90!black, font=\bfseries\scriptsize, inner sep=4pt}
  ]
    \node[b1] (map) at (-4, 0) {BẢN ĐỒ (BÉ)};
    \node[b2] (real) at (4, 0) {THỰC TẾ (LỚN)};

    \draw[teal!85!black, very thick, ->] (map) to[bend left=25] node[above, font=\scriptsize\bfseries] {$\times n$ (Muốn ra thực tế thì Nhân)} (real);
    \draw[teal!85!black, very thick, ->] (real) to[bend left=25] node[below, font=\scriptsize\bfseries] {$: n$ (Muốn vẽ lên bản đồ thì Chia)} (map);
  \end{tikzpicture}
\end{tcolorbox}

\vspace{3mm}

% --- KHỐI 4: VÍ DỤ MINH HỌA & BẪY TRẮC NGHIỆM ---
\noindent
\begin{minipage}[t]{0.54\textwidth}
  \begin{tcolorbox}[
    enhanced,
    title={\footnotesize\bfseries\color{white}\faLightbulbO\ Ví dụ mẫu giải chi tiết},
    colback=green!3!white,
    colframe=green!60!black,
    arc=2.5mm, boxrule=0.8pt, top=1.8mm, bottom=1.8mm, left=2.2mm, right=2.2mm
  ]
    \footnotesize
    Trên bản đồ tỉ lệ $1 : 2\,000\,000$, quãng đường từ Hà Nội đến Hải Phòng đo được $5\text{ cm}$. Hỏi ngoài thực tế dài bao nhiêu km?\\
    $\triangleright$ Quãng đường thực tế là: $5 \times 2\,000\,000 = 10\,000\,000\text{ cm}$.\\
    $\triangleright$ Đổi: $10\,000\,000\text{ cm} = 100\,000\text{ m} = \mathbf{100\text{ km}}$.
  \end{tcolorbox}
\end{minipage}%
\hfill
\begin{minipage}[t]{0.44\textwidth}
  \begin{tcolorbox}[
    enhanced,
    title={\footnotesize\bfseries\color{white}\faExclamationTriangle\ Mẹo tránh bẫy trắc nghiệm},
    colback=red!2!white,
    colframe=red!70!black,
    arc=2.5mm, boxrule=0.8pt, top=1.8mm, bottom=1.8mm, left=2.2mm, right=2.2mm
  ]
    \footnotesize
    $\bullet$ Đếm thiếu/thừa chữ số $0$ khi đổi từ $\text{cm}$ sang $\text{km}$ ($1\text{ km} = 100\,000\text{ cm}$, bỏ $5$ chữ số $0$).\\
    $\bullet$ Khi tính độ dài trên bản đồ: Quên đổi độ dài thực tế ra $\text{cm}$ trước khi chia cho $n$.
  \end{tcolorbox}
\end{minipage}

\end{document}
